\documentclass[10pt,a4paper,english]{article}

\usepackage[margin=2cm]{geometry}
\usepackage[T1]{fontenc}
\usepackage[utf8]{inputenc}
\usepackage[numbers,sort&compress]{natbib}
\usepackage{babel}
\usepackage{lmodern}
\usepackage{microtype}

\usepackage{amsmath,amssymb}
\usepackage{graphicx}
\usepackage{booktabs}
\usepackage{subcaption}
\usepackage{xcolor}
\usepackage[colorlinks=true,linkcolor=blue!50!black,citecolor=blue!50!black,urlcolor=blue!50!black]{hyperref}
\usepackage[capitalize,nameinlink]{cleveref}

\graphicspath{{figures/}}
\setlength{\parskip}{0.7em}

\newcommand{\E}{\mathbb{E}}
\newcommand{\D}{\mathcal{D}}
\newcommand{\piexp}{\pi^{*}}
\newcommand{\ind}{\mathbf{1}}

\title{DAgger vs symmetric sampling in Super Tux Kart}
\author{Théo Fontaine \and Théophile Lahaussois\\[0.3em]
  \small Sorbonne Université -- M2 MIND -- Reinforcement Learning\\
  \small Code: \url{https://github.com/theolahaussois/DAgger-vs-symmetric-sampling-in-Super-Tux-Kart}}
\date{}

\begin{document}
\maketitle

\begin{abstract}
We compare DAgger and the \emph{symmetric sampling} of RLPD for teaching a single SuperTuxKart kart to complete one lap as fast as possible, from the same hand-written expert and the same demonstration lap. Over 10 seeds per method, DAgger matches the expert from its first iteration and, with a decaying mixing coefficient, finishes every lap and obtains the highest and most reliable scores at equal budget. RLPD needs about 8\,600 steps to finish a first lap and then catches up with DAgger in median, but two seeds out of ten collapse. Without the demonstrations, the same algorithm never finishes a lap. Both methods sometimes beat the expert by almost 3\,s by avoiding a banana that the expert hits on every lap.
\end{abstract}

\section{Introduction}
\label{sec-intro}

To learn to drive from an expert $\piexp$, one can \emph{imitate} it by asking which action it would take in the states we visit, as DAgger does \citep{ross2011dagger}, or \emph{learn by reinforcement} while reusing its trajectories as past experience, as the \emph{symmetric sampling} of RLPD does \citep{ball2023rlpd}. DAgger queries the expert online and ignores the reward; RLPD never queries it but directly optimizes the reward. We compare them on a single SuperTuxKart kart that must complete one lap as fast as possible, with four questions: \textbf{(Q1)} for an equal interaction budget, which method produces the fastest driver? \textbf{(Q2)} what do we gain by querying the expert online rather than imitating its demonstrations, and what is the effect of DAgger's mixing coefficient $\beta_i$? \textbf{(Q3)} what does RLPD owe to symmetric sampling? \textbf{(Q4)} can a method outperform the expert?

\section{Methods}
\label{sec-methods}

Behavioral cloning (BC) learns a deterministic policy $\pi_\theta$ by regression on the expert's pairs $(s_t, a_t)$, minimizing $\mathcal{L}(\theta) = \sum_t \|a_t - \pi_\theta(s_t)\|^2$. It is trained on the expert's states but deployed on its own: a small error leads it to out-of-distribution (OOD) states, where it may act poorly and drift further away. This \emph{policy-related distribution shift} makes the gap to the expert grow as $O(H^2\varepsilon)$ over a horizon $H$ \citep{ross2010efficient}. DAgger corrects it. At iteration $i$, we drive with the mixture $\pi_i = \beta_i \piexp + (1-\beta_i)\hat{\pi}_i$ (the expert acts with probability $\beta_i$), label each visited state with the action $\piexp(s)$, add it to $\D$, then retrain $\hat{\pi}_{i+1}$ on the whole of $\D$. The learner thus learns on its own state distribution, and the gap becomes linear in $H$. We use $\beta_i = p^{\,i}$ and compare $p = 0$, where the expert no longer drives after the demonstrations (often the best choice according to \citet{ross2011dagger}), with $p = 0.5$.

RLPD is a SAC \citep{haarnoja2018sacapps} that is given a fixed set of transitions $\D_{\text{off}}$ in addition to the online buffer $\D_{\text{on}}$. Symmetric sampling builds each batch of size $b$ from $b/2$ transitions of $\D_{\text{off}}$ and $b/2$ of $\D_{\text{on}}$. If the demonstrations were simply poured into the buffer, their share, $|\D_{\text{off}}|/(|\D_{\text{off}}|+|\D_{\text{on}}|)$, would drop below $2\%$ by the end of training here. This mixing is possible because the critic's target does not depend on the behavioral policy $\pi_\beta$ that generated the transition: RLPD does not imitate the expert, it learns the value of its actions. RLPD adds Layer Normalization to the critics, which limits the \emph{critic-related distribution shift} (Q-values of OOD actions stay bounded instead of being overestimated), and a high update-to-data ratio (UTD, here the number of critic updates per environment step). To fit our compute budget, we use two critics instead of an ensemble of ten, and a UTD of 10 instead of 20. An ablation, \emph{SAC only}, removes the demonstrations.

DAgger receives an action label at every state: a rich signal, but one that only says what the expert would do, and that costs one expert query per step. RLPD only receives a scalar reward, propagated by the Bellman equation: a poorer signal, but one that says what is good. DAgger is therefore capped by the expert, whereas RLPD can surpass it, at the cost of exploration and a critic that are harder to control.

\section{Experimental protocol}
\label{sec-protocol}

We use \texttt{pystk2-gymnasium} \citep{pystk2gym} (\texttt{supertuxkart/full-v0}) on the track \texttt{cocoa\_temple}, with one kart, one lap and at most 1500 steps of $0.1$\,s. With a single kart, the native reward is $r_{t+1} = (\delta_{t+1} - \delta_t)/10 - 0.1 + 10 \cdot \ind[\text{finish}]$, where $\delta_t$ is the distance travelled. The return of a lap of length $L \approx 1517$\,m completed in $T$ steps is therefore
\begin{equation}
  G = L/10 - 0.1\,T + 10 ,
  \label{eq-return}
\end{equation}
and that of an unfinished lap is $G = \delta/10 - 150$: for a normal lap, maximizing $G$ is equivalent to minimizing the lap time. The environment is deterministic: a given deterministic policy always drives the same lap (only graphical rendering, never used during evaluation, alters the simulation).

The learner receives 28 numbers expressed in the kart's frame: velocity, vector and distance to the track centre, maximum steering angle, endpoints and widths of the next 5 track segments. It acts with $a = (\text{steer}, \text{throttle}) \in [-1,1]^2$, a negative throttle meaning braking; drift, nitro, firing and rescue are disabled. The expert is \texttt{Agent1} (team 1 of \texttt{pystk\_project} \citep{sigaud2025pystk}), a hand-written controller, from which we remove item and nitro handling and whose actions we project onto the learner's action space. The learner can thus closely reproduce the expert, as DAgger's guarantee requires; only approximately, though, since the expert keeps a small internal memory and reads observations that the learner does not receive.

The demonstration data is one lap of the expert, i.e. 609 transitions. Since the environment is deterministic, further laps would be copies. We do not add noise to the expert's actions: this technique (DART \citep{laskey2017dart}) itself corrects BC's distribution shift, and would bias the comparison with DAgger. This dataset serves as the initial $\D$ for BC and DAgger, and as $\D_{\text{off}}$ for RLPD.

Every policy is evaluated on one lap, without noise (mean action for SAC). Determinism makes repeating this episode unnecessary: all variability comes from training. The main metric is the return \eqref{eq-return}, complemented by the share of finished laps and the lap time. Each method is trained with 10 seeds, which are the statistical unit; a seed's score is that of its final policy, without selecting the best one. The common axis is the number of online environment steps: about 12\,000 for DAgger and 30\,000 for RLPD, whose score we also read at 12\,000 steps to compare at equal budget. We report the mean with a 95\,\% bootstrap confidence interval, the median and the IQM \citep{agarwal2021rliable}. Methods are compared with Welch and Mann--Whitney tests \citep{colas2019hitchhiker}, and with the (deterministic) expert with a one-sample Student test and a Wilcoxon test. All comparisons, fixed in advance, are corrected with the Holm method \citep{holm1979simple}.

The code is written in PyTorch (no RL library, rather than the recommended rl-mind toolkit, so that the three methods share the same minimal code) and trained on the CPU of an Apple M1; hyper-parameters are given in \cref{tab-hparams}.

\begin{table}[h]
  \centering
  \small
  \caption{Hyper-parameters. Those of SAC are the usual values, not tuned.}
  \label{tab-hparams}
  \begin{tabular}{lll}
    \toprule
    Common & networks & 2 hidden layers of 256, ReLU; inputs standardized on the demonstrations \\
    BC and DAgger & training & squared error, $\tanh$ output, Adam ($10^{-3}$), batch 256, 2000 steps \\
     & DAgger & 20 iterations of one episode, full retraining, $p \in \{0, 0.5\}$ \\
    SAC and RLPD & optimization & Adam ($3 \cdot 10^{-4}$), $\gamma = 0.99$, $\tau = 0.005$, batch 256 (RLPD: 128 + 128) \\
     & critics & 2, minimum in the target, Layer Normalization, 10 updates per step \\
     & actor & $\tanh$-squashed Gaussian, $\alpha_0 = 1$, target entropy $-2$ \\
     & budget & 30\,000 steps, evaluation about every 1000 steps \\
    \bottomrule
  \end{tabular}
\end{table}

\section{Results}
\label{sec-results}

% Figures and tables are produced by analysis/results.py and analysis/trajectory_gap.py.
\newcommand{\fig}[2]{\IfFileExists{figures/#1.pdf}{\includegraphics[width=#2]{#1.pdf}}{\fbox{\parbox[c][3.5cm][c]{#2}{\centering\ttfamily\detokenize{#1}.pdf}}}}

\Cref{tab-summary} summarizes the final scores, and \cref{fig-curves} shows how each method gets there as a function of environment steps. The expert is the reference: a return of 100.83, i.e. 60.9\,s.

\begin{table}[t]
  \centering
  \caption{Final scores over 10 seeds (one deterministic evaluation per seed). Return: mean and 95\,\% bootstrap CI. Success: seeds that finish the lap. Time: mean over finished laps. 1\textsuperscript{st} lap: median number of online steps before the first finished lap. Last 5: mean of the last 5 evaluations of each seed (exploratory analysis). RLPD @12k: RLPD read at DAgger's budget.}
  \label{tab-summary}
  \resizebox{\linewidth}{!}{\IfFileExists{figures/summary_table.tex}{\begin{tabular}{lccccccc}
\toprule
Method & Return [95\,\% CI] & Median & IQM & Success & Time (s) & 1\textsuperscript{st} lap (steps) & Last 5 [95\,\% CI] \\
\midrule
Expert & 100.83 & -- & -- & 100\,\% & 60.9 & -- & -- \\
BC & 101.20 [101.03, 101.40] & 101.16 & 101.17 & 100\,\% & 60.5 & 0 & -- \\
DAgger p=0 & 90.67 [68.38, 102.22] & 101.64 & 101.45 & 90\,\% & 60.0 & 0 & 99.81 [95.37, 102.29] \\
DAgger p=0.5 & 102.13 [101.02, 103.11] & 102.88 & 102.51 & 100\,\% & 59.6 & 0 & 102.11 [101.57, 102.59] \\
RLPD @12k & 65.72 [18.58, 100.93] & 100.98 & 98.70 & 80\,\% & 62.2 & -- & -- \\
RLPD & 51.84 [$-$23.55, 102.76] & 102.71 & 101.84 & 80\,\% & 59.4 & 8596 & 56.91 [$-$8.53, 102.45] \\
SAC only & $-$142.22 [$-$149.93, $-$130.25] & $-$149.97 & $-$148.70 & 0\,\% & -- & -- (10 never) & $-$139.88 [$-$147.14, $-$131.57] \\
\bottomrule
\end{tabular}
}{summary\_table.tex}}
\end{table}

\begin{figure}[t]
  \centering
  \begin{subfigure}[t]{0.53\linewidth}\centering\fig{learning_curves}{\linewidth}\caption{Return during training. Top: mean and 95\,\% CI. Bottom: median, zoomed around the expert (dashed).}\label{fig-learning}\end{subfigure}\hfill
  \begin{subfigure}[t]{0.45\linewidth}\centering\fig{success_rate}{\linewidth}\caption{Share of seeds whose current policy finishes the lap (mean and 95\,\% CI).}\label{fig-success}\end{subfigure}
  \caption{Learning curves as a function of the number of online environment steps (10 seeds per method). The vertical dotted line marks DAgger's budget. Both DAgger variants stay at the expert's level from the start; RLPD only learns to drive after several thousand steps, then remains unstable; SAC only never succeeds.}
  \label{fig-curves}
\end{figure}

BC, which is DAgger's iteration 0, already finishes all its laps and slightly outperforms the expert (101.20 versus 100.83). The following DAgger iterations stay at this level (\cref{fig-learning}) and finish 209 laps out of 210 evaluations for $p = 0$, and 210 out of 210 for $p = 0.5$. With $p = 0.5$, DAgger obtains the best scores of the study: mean 102.13, median 102.88, IQM 102.51, i.e. about 1.3\,s faster than the expert. The $p = 0$ variant has a similar median (101.64), but its last iteration failed on one seed (kart stuck 100\,m from the finish line), which drops its mean to 90.67. From one iteration to the next, the lap time varies between 58 and 66\,s: each retraining starts from a fresh network, and a small steering difference can change the whole end of the lap.

RLPD starts with a random actor and stays still for several thousand steps, presumably the time needed for the temperature $\alpha$ to decrease and for the critic to learn the value of the expert's actions (we did not log these quantities). All seeds then complete a first lap, between 6\,800 and 12\,500 steps (median 8\,600). An unstable phase follows, visible on the median between 12\,000 and 20\,000 steps, where policies alternate between fast laps and failures (\cref{fig-success}). At the end of training, 8 seeds out of 10 finish the lap, in 59.4\,s on average, and 5 of them in under 58.5\,s; the other two have collapsed and stay still. Hence a high median (102.71), close to DAgger, but a mean of 51.84.

With the same algorithm and the same budget, but without expert data, no seed finishes a single lap: the best kart covers 1\,316\,m out of 1\,517, and most end up still (final return $-142.22$).

\begin{figure}[t]
  \centering
  \begin{subfigure}[t]{0.49\linewidth}\centering\fig{final_scores}{\linewidth}\caption{Final score of each seed (dots) and mean with 95\,\% CI (bars). Bottom: zoom around the expert.}\label{fig-final}\end{subfigure}\hfill
  \begin{subfigure}[t]{0.49\linewidth}\centering\fig{trajectory_gap_rlpd}{\linewidth}\caption{RLPD (seed 0, 58.0\,s) versus the expert, as a function of the distance travelled: time gap, steering and throttle. The grey area is where the expert gets the parachute.}\label{fig-gap}\end{subfigure}
  \caption{Left: the spread of final scores reveals RLPD's collapses and the single failure of DAgger $p=0$. Right: RLPD's lead over the expert builds up mostly at the banana.}
  \label{fig-final-gap}
\end{figure}

The expert drives over the same banana on every lap and gets a parachute for 3.8\,s (38 steps). The two fast policies we inspected (DAgger $p = 0.5$, seed 4, 58.3\,s, and RLPD, seed 0, 58.0\,s) take another line at this spot and never hit the banana; the many final laps clustered around 58.2\,s suggest the same for the others, but we did not check them all. \Cref{fig-gap} shows it for RLPD: the time gap stays small (about 0.2\,s) for 1\,300\,m, then widens to about 2.9\,s precisely in this area. The BC policy we inspected (seed 0) gets the parachute like the expert, and its advantage comes from the end of the lap, presumably because the network smooths the expert's threshold-based rules.

After Holm correction over the 14 comparisons defined in advance (\cref{tab-tests}), only SAC only is significantly worse than RLPD ($-194$, $p = 0.003$) and than the expert. BC beats the expert by 0.37 with a CI excluding 0, but stays at the edge of the threshold ($p = 0.05$ to $0.06$). The other differences are not significant: RLPD's collapses make the intervals very wide, and 10 seeds only detect effects of about 1.25 standard deviations.

\begin{table}[t]
  \centering
  \footnotesize
  \caption{Comparisons between methods and with the expert: difference of means ($\Delta$, with 95\,\% bootstrap CI) and $p$-values corrected with the Holm method over all 14 tests. Against the deterministic expert, the tests are one-sample (Student and Wilcoxon).}
  \label{tab-tests}
  \IfFileExists{figures/tests_table.tex}{\begin{tabular}{lccc}
\toprule
Comparaison & $\Delta$ [IC 95\,\%] & $p$ Welch (Holm) & $p$ Mann--Whitney (Holm) \\
\midrule
RLPD $-$ DAgger p=0.5 & $-$50{,}29 [$-$125{,}51 ; 0{,}76] & 1{,}000 & 1{,}000 \\
RLPD $-$ DAgger p=0 & $-$38{,}83 [$-$113{,}71 ; 21{,}96] & 1{,}000 & 1{,}000 \\
RLPD @12k $-$ DAgger p=0.5 & $-$36{,}41 [$-$83{,}47 ; $-$1{,}34] & 1{,}000 & 0{,}190 \\
RLPD @12k $-$ DAgger p=0 & $-$24{,}95 [$-$80{,}62 ; 20{,}30] & 1{,}000 & 1{,}000 \\
SAC seul $-$ RLPD & $-$194{,}06 [$-$249{,}47 ; $-$122{,}38] & 0{,}003 & 0{,}092 \\
DAgger p=0.5 $-$ DAgger p=0 & 11{,}46 [$-$0{,}55 ; 34{,}06] & 1{,}000 & 1{,}000 \\
DAgger p=0 $-$ BC & $-$10{,}53 [$-$32{,}88 ; 1{,}03] & 1{,}000 & 1{,}000 \\
DAgger p=0.5 $-$ BC & 0{,}93 [$-$0{,}20 ; 1{,}94] & 1{,}000 & 1{,}000 \\
DAgger p=0.5 $-$ Expert & 1{,}30 [0{,}20 ; 2{,}28] & 0{,}520 & 0{,}488 \\
BC $-$ Expert & 0{,}37 [0{,}19 ; 0{,}57] & 0{,}062 & 0{,}051 \\
RLPD $-$ Expert & $-$48{,}99 [$-$124{,}36 ; 1{,}92] & 1{,}000 & 1{,}000 \\
DAgger p=0 $-$ Expert & $-$10{,}16 [$-$32{,}45 ; 1{,}40] & 1{,}000 & 1{,}000 \\
SAC seul $-$ Expert & $-$243{,}05 [$-$250{,}76 ; $-$231{,}08] & 0{,}000 & 0{,}027 \\
RLPD @12k $-$ Expert & $-$35{,}11 [$-$82{,}72 ; 0{,}05] & 1{,}000 & 1{,}000 \\
\bottomrule
\end{tabular}
}{tests\_table.tex}
\end{table}

\section{Discussion}
\label{sec-discussion}

\paragraph{Q1.} After about 12\,000 steps, DAgger finishes all its laps, whereas RLPD only finishes 8 out of 10 and has an IQM almost 4 points lower (98.70 versus 102.51). The confidence interval of the difference of means excludes 0 ($-36.41$ [$-83.47$, $-1.34$]), even though the test does not survive the Holm correction. DAgger is also much cheaper in compute: a few minutes per seed, versus about 40 minutes for RLPD with 10 critic updates per step. In return, it requires an expert that can be queried at every step, i.e. about 12\,600 labels, whereas RLPD makes do with a single recorded lap. With 2.5 times more steps, RLPD catches up with DAgger in median, but not in reliability.

\paragraph{Q2.} DAgger does not do significantly better than BC ($+0.93$, CI [$-0.20$, $1.94$]). Theory explains why: BC's $O(H^2\varepsilon)$ bound describes the worst case, where every error pushes the learner away from the expert's trajectory. Here, the environment is deterministic, the expert mostly reacts to the current state and the learner reproduces it closely: imitation is accurate enough that the learner never leaves the trajectory, and the policy-related distribution shift does not appear. The choice of $\beta_i$ has no demonstrable effect: a single failed seed separates the two variants.

\paragraph{Q3.} This is the clearest effect of the study. The expert's 609 transitions are enough to take SAC from zero finished laps to ten seeds out of ten that finish at least one. Since they make up half of every batch, the critic learns from the start the value of a trajectory that reaches the finish line, which random exploration does not discover in 30\,000 steps.

\paragraph{Q4.} RLPD does so as expected, by directly optimizing the lap time. The case of DAgger is more surprising: DAgger's guarantee bears on the imitation error, not on performance, and nothing prevents the learner from doing better than its teacher. By distilling a rule-based controller into a neural network, DAgger plausibly obtains a smoother driving (not measured here), which sometimes finds a line that avoids the banana.

Our RLPD is simplified (2 critics instead of 10, UTD of 10 instead of 20). For SAC alone, the best UTD is usually 1 to 3, because a higher ratio increases overestimation and overfitting \citep{li2023overfitting}; the critic ensemble is what makes higher ratios usable. Our two collapses are therefore probably linked to this choice, and testing the full RLPD is the natural next step. The task also limits the comparison: BC already solves it, so DAgger's correction of the distribution shift is barely exercised; demonstrations covering only part of the lap, or perturbed dynamics, would test it more directly. The study covers one track and one expert. DAgger's last-iteration score is noisy (the mean of the last 5 evaluations is more stable: 102.11 for $p = 0.5$). Finally, two intermediate RLPD evaluations measured more distance than the track length, which inflates their return; no final policy is affected.

\section{Conclusion}
\label{sec-conclusion}

On this deterministic task with a reactive expert, DAgger is the most efficient and reliable method, provided the expert can be queried at every step. RLPD only needs one recorded lap and can also beat the expert, but it needs more interactions and compute, and remains unstable in our simplified version. Symmetric sampling is what makes it work: without it, the same algorithm never completes a lap. Since DAgger can be seen as a way to distill a hand-written controller into a neural network, a natural extension is to fine-tune the DAgger policy with RL, to combine a reliable start with reward optimization.

\renewcommand{\bibfont}{\small}
\setlength{\bibsep}{1pt}
\bibliographystyle{plainnat}
\bibliography{refs}

\end{document}
